\ChapterImagePrelim[cap:introduccion]{Introducción}{./images/fondo.png}\label{cap:introduccion}
\mbox{}\\

La infraestructura de \TI~constituye uno de los pilares esenciales del entorno digital actual, ya que sustenta el funcionamiento de sistemas y servicios de los que dependen las organizaciones. En un sentido amplio, la infraestructura abarca todos los elementos técnicos —Hardware, Software, Redes, Instalaciones, etc.— necesarios para el desarrollo, prueba, entrega, supervision, control o soporte para aplicaciones y servicios de \TI~\cite{Axelos01}~. Según lo anterior se puede evidencia la importancia de una adecuada planificación y gestión para la eficiencia, seguridad y sostenibilidad de los sistemas de información, lo que en consecuencia configura a la infraestructura de \TI~como una actividad estratégica tanto en el ámbito empresarial como en el académico.

Sin embargo, a pesar de los avances en metodologías y marcos de referencia, el diseño de la infraestructura tecnológica continúa enfrentando desafíos relacionados con la representación formal y la comunicación de sus componentes. Diversos autores han señalado la falta de un lenguaje estandarizado que unifique criterios de modelado, generando inconsistencias y ambigüedades que dificultan la interoperabilidad conceptual entre herramientas y profesionales \cite{Escobar01, Parviainen01}~. Esta problemática repercute directamente en la documentación técnica, la comprensión de los modelos y la transferencia de conocimiento dentro de las comunidades de práctica y aprendizaje.

Frente a este escenario, el presente trabajo de grado, titulado Notación para el Diseño de Infraestructura de TI (\NDITI): una propuesta desde la Universidad del Quindío, plantea el desarrollo de una notación de modelado unificada orientada a resolver las limitaciones de interoperabilidad y uniformidad conceptual detectadas en las notaciones existentes. Para ello, se plantea una metodología que combina enfoques tradicionales de investigación con prácticas ágiles, articulando tres ejes principales: el reconocimiento de necesidades del entorno académico, la realización de un estudio de mapeo sistemático (\SMS) para identificar referentes metodológicos y tecnológicos, y la aplicación de la técnica Decision Analysis and Resolution (\DAR) del modelo CMMI para seleccionar y adaptar la notación más adecuada.

El resultado de este proceso se materializa en la propuesta de una notación denominada NDITI, diseñada para servir como lenguaje formal de modelado en el contexto académico de la Universidad del Quindío. Esta notación pretende integrar los beneficios de lenguajes previos adaptándolos a un enfoque pedagógico que priorice la comprensión y la coherencia entre los distintos niveles de abstracción del diseño. Asimismo, se busca que la NDITI sirva como instrumento de comunicación entre docentes, estudiantes y profesionales, fortaleciendo las competencias de análisis y diseño en el campo de la infraestructura tecnológica.